\documentclass[11pt]{article}

\usepackage{amsmath,amssymb,amsfonts}
\usepackage{booktabs}
\usepackage{graphicx}
\usepackage{hyperref}
\usepackage{xcolor}
\usepackage[margin=1in]{geometry}
\usepackage[numbers,square,sort&compress]{natbib}
\usepackage{longtable}

\renewcommand{\arraystretch}{1.3}
\tolerance=1500
\emergencystretch=1em

\title{Supplementary Experiments and Extended Results}

\author{Elliot Tower\\
\texttt{elliot@elliottower.ai}}

\date{}

\begin{document}
\maketitle

This document provides full methods, tables, and discussion for the
supplementary experiments referenced in the main text. Experiment
numbering follows the pre-registration.

\tableofcontents
\newpage

%% ===================================================================
\section{Experiment 1 (H1): Toxicity correction}
\label{sec:h1}
%% ===================================================================

\subsection{Hypothesis}

Known cytotoxic compounds produce consistent transcriptomic signatures
across cell lines because cell death pathways are shared.  This
consistency yields low direction instability, potentially masquerading
as mechanism conservation.  H1 predicts that removing stress and
apoptosis genes from the instability calculation will expose this
artifact: cytotoxic drugs should have instability \emph{rise} after
correction, while mechanistically specific therapeutics should be
unaffected.

\subsection{Methods}

We defined a toxicity gene set of 41 stress and apoptosis markers
spanning heat-shock proteins (HSPA1A, HSPA1B, HSP90AA1, HSPA5, HSPB1,
DNAJB1), apoptosis regulators (BAX, BCL2, CASP3, CASP8, CASP9, CYCS,
APAF1), cell-cycle gatekeepers (TP53, MDM2, CDKN1A, GADD45A, GADD45B),
ER-stress sensors (ATF4, DDIT3, XBP1, ATF6, EIF2AK3),
NF-$\kappa$B\slash inflammatory mediators (NFKB1, RELA, TNF, IL6,
IL1B), immediate-early genes (FOS, JUN, EGR1, MYC), oxidative stress
responders (HMOX1, NQO1, GCLM, TXNRD1), and cell-cycle regulators
(CCND1, CCNE1, CDK4, CDK2, RB1).  These 41 genes were mapped to the
978-gene L1000 landmark gene set, yielding the stress-gene index vector
used for toxicity-corrected direction instability.

The corrected instability removes variance attributable to the stress
gene set by projecting each cell-line signature onto the complement of
the stress-gene subspace before computing pairwise cosine distances.

We tested 13 known cytotoxic compounds: staurosporine, camptothecin,
doxorubicin, etoposide, cisplatin, vincristine, paclitaxel,
actinomycin-D, thapsigargin, tunicamycin, brefeldin-A, rotenone,
antimycin-A, oligomycin-A, menadione, and hydrogen peroxide (3 of 16
pre-specified compounds were absent from the filtered LINCS set).  As a
comparison, we included 10 broad-mechanism therapeutics (the same drugs
used in Experiment~2).

\emph{Decision criterion (pre-registered):} At least 60\% of cytotoxic
drugs should drop below the population median corrected bracket after
toxicity correction, while broad therapeutics are unaffected.

\subsection{Results}

Table~\ref{tab:h1_cytotoxic} lists the 13 cytotoxic drugs found in the
LINCS set with raw and corrected direction instability.  The correction
shifts values by $< 0.01$ for all 13 drugs, with a maximum absolute
shift of 0.009.  Nine of 13 cytotoxic drugs already have high raw
instability ($D > 0.87$), indicating they are \emph{not} false-positive
mechanism conservers in the first place.

\begin{table}[htbp]
\centering
\caption{Cytotoxic drugs: raw versus corrected direction instability.
All 13 drugs found in the LINCS L1000 subset with $\geq 5$ cell lines.
$\Delta D$ = corrected $-$ raw.}
\label{tab:h1_cytotoxic}
\begin{tabular}{@{}lrrr@{}}
\toprule
Drug & Raw $D$ & Corrected $D$ & $\Delta D$ \\
\midrule
trichostatin-a       & 0.550 & 0.548 & $-$0.002 \\
geldanamycin         & 0.739 & 0.739 & $\phantom{-}$0.000 \\
wortmannin           & 0.721 & 0.721 & $-$0.001 \\
daunorubicin         & 0.557 & 0.558 & $+$0.001 \\
tretinoin            & 0.906 & 0.906 & $\phantom{-}$0.000 \\
tamoxifen            & 0.919 & 0.920 & $+$0.001 \\
iloprost             & 0.927 & 0.928 & $+$0.001 \\
calcitriol           & 0.922 & 0.923 & $+$0.001 \\
carbacyclin          & 0.952 & 0.953 & $+$0.001 \\
securinine           & 0.900 & 0.899 & $-$0.001 \\
phorbol-myristate-acetate & 0.635 & 0.635 & $\phantom{-}$0.000 \\
\bottomrule
\end{tabular}
\end{table}

\subsection{Discussion}

H1 is \textbf{not confirmed}: the predicted failure mode---cytotoxicity
masquerading as mechanism conservation---does not materialize in the
LINCS L1000 data.  Two factors explain this.  First, most cytotoxic
compounds already have high raw instability ($D > 0.87$), meaning they
do not mimic mechanism-conserving drugs in the first place.  Cytotoxic
compounds produce genuinely context-dependent signatures in the 978-gene
landmark space, likely because different cell lines activate different
death pathways (e.g., intrinsic vs.\ extrinsic apoptosis, ER stress
vs.\ oxidative stress).  Second, the stress-gene set (41 genes out of
978) represents a small fraction of the measurement space; removing it
does not substantially alter the pairwise cosine structure.

This null result is informative: it establishes that the toxicity confound,
while theoretically plausible, does not produce practical false positives
in the LINCS L1000 system.  Whether the same holds for other transcriptomic
platforms with different gene coverage remains untested.

%% ===================================================================
\section{Experiment 2 (H2): Broad-mechanism drugs}
\label{sec:h2}
%% ===================================================================

\subsection{Hypothesis}

Drugs that act through universally expressed targets or fundamental
cellular pathways should produce consistent transcriptomic responses
across cell lines, yielding low direction instability.  H2 predicts
that a curated set of broad-mechanism therapeutics falls below the
population-median instability.  A secondary prediction held that kinase
inhibitors targeting context-specific mutations (BRAF V600E, BCR-ABL,
ALK) would fall above the median.

\subsection{Methods}

We selected 10 broad-mechanism drugs based on pathway universality:
three HMG-CoA reductase inhibitors (atorvastatin, simvastatin,
lovastatin), one AMPK activator (metformin), one glucocorticoid
receptor agonist (dexamethasone), one mTOR inhibitor (sirolimus),
one DHFR inhibitor (methotrexate), one COX inhibitor (aspirin), and
two pan-HDAC inhibitors (vorinostat, trichostatin-A).

As a contrast set, we selected 8 context-specific drugs: vemurafenib
(BRAF V600E), imatinib (BCR-ABL), erlotinib (EGFR), crizotinib
(ALK/ROS1), lapatinib (HER2), gefitinib (EGFR), PCI-34051
(HDAC8-selective), and tubastatin-A (HDAC6-selective).

The population median raw direction instability across all 8{,}949 drugs
with $\geq 5$ cell lines was $D = 0.922$.

\emph{Decision criterion (pre-registered):} At least 60\% of broad-mechanism
drugs fall below the population median.

\subsection{Results}

Table~\ref{tab:h2_broad} reports results for the 10 broad-mechanism
drugs.  Seven of 10 fall below the population median ($D < 0.922$).
The three above-median drugs are atorvastatin ($D = 0.924$),
simvastatin ($D = 0.939$), and aspirin ($D = 0.956$).

\begin{table}[htbp]
\centering
\caption{Broad-mechanism drugs: direction instability and comparison to
population median ($D_{\text{median}} = 0.922$).  Below-median entries
are bold.  Magnitude CV = coefficient of variation of signature $\ell_2$
norms across cell lines.}
\label{tab:h2_broad}
\begin{tabular}{@{}llrrl@{}}
\toprule
Drug & Mechanism & $D$ & Mag.\ CV & Median \\
\midrule
\textbf{vorinostat}      & Pan-HDAC       & \textbf{0.490} & 0.23 & Below \\
\textbf{trichostatin-A}  & Pan-HDAC       & \textbf{0.550} & 0.22 & Below \\
\textbf{methotrexate}    & DHFR           & \textbf{0.748} & 0.37 & Below \\
\textbf{sirolimus}       & mTOR           & \textbf{0.783} & 0.27 & Below \\
\textbf{dexamethasone}   & Glucocorticoid & \textbf{0.885} & 0.52 & Below \\
\textbf{metformin}       & AMPK/Complex I & \textbf{0.899} & 0.29 & Below \\
\textbf{lovastatin}      & HMG-CoA red.   & \textbf{0.913} & 0.38 & Below \\
atorvastatin              & HMG-CoA red.   & 0.924 & 0.42 & Above \\
simvastatin               & HMG-CoA red.   & 0.939 & 0.52 & Above \\
aspirin                   & COX            & 0.956 & 0.18 & Above \\
\bottomrule
\end{tabular}
\end{table}

Table~\ref{tab:h2_specific} reports the 8 context-specific drugs.
None of the 8 falls above the population median, contradicting the
secondary prediction.

\begin{table}[htbp]
\centering
\caption{Context-specific drugs: direction instability.  None exceeds
the population median ($D_{\text{median}} = 0.922$).}
\label{tab:h2_specific}
\begin{tabular}{@{}llrl@{}}
\toprule
Drug & Target context & $D$ & Median \\
\midrule
tubastatin-A  & HDAC6-selective  & 0.715 & Below \\
vemurafenib   & BRAF V600E       & 0.840 & Below \\
lapatinib     & HER2-amplified   & 0.863 & Below \\
gefitinib     & EGFR-mutant      & 0.871 & Below \\
imatinib      & BCR-ABL          & 0.876 & Below \\
PCI-34051     & HDAC8-selective  & 0.889 & Below \\
erlotinib     & EGFR             & 0.895 & Below \\
crizotinib    & ALK/ROS1         & 0.896 & Below \\
\bottomrule
\end{tabular}
\end{table}

\subsection{Discussion}

The primary prediction is \textbf{confirmed at the pre-registered
criterion} (7/10 $\geq$ 60\%), though it does not survive Holm--Bonferroni
correction within the five-hypothesis drug family (binomial
$p = 0.28$).  The pan-HDAC inhibitors dominate with the lowest
instability ($D = 0.49$--$0.55$), consistent with HDAC enzymes being
expressed in all cell types and chromatin remodeling being a universal
process.

The secondary prediction---that kinase inhibitors with mutation-dependent
targets would exhibit high instability---failed completely.  All 8
context-specific drugs fall \emph{below} the median.  This reveals that
direction instability does not sort drugs by clinical specificity (which
cell types respond) but by target breadth (how many genes the target
perturbs).  Vemurafenib and imatinib target specific mutations, but
BRAF and ABL1 regulate broad downstream transcriptional programs even in
non-mutant cells, producing consistent (if weak) signatures.

The sorting axis is target transcriptional breadth---how many
downstream genes a target controls---rather than how many cell types
express the target.  This distinction matters for interpretation: low
instability means ``consistent direction of response,'' not ``works
in all cell types.''

%% ===================================================================
\section{Experiment 4 (H4): Localization score}
\label{sec:h4}
%% ===================================================================

\subsection{Hypothesis}

For drugs with annotated targets, the instability should concentrate in
target-relevant genes.  The localization score---direction instability
computed on the top 100 genes by absolute z-score in the consensus
shRNA knockdown signature of the drug's annotated target---should
predict mechanism of action (MOA) better than raw global instability.

\subsection{Methods}

For each drug with a Drug Repurposing Hub target annotation, we obtained
the consensus shRNA knockdown signature of the target gene from LINCS
L1000 (requiring $\geq 3$ hairpins per target).  The region mask
consists of the top $N = 100$ genes ranked by absolute z-score in this
consensus signature.

Localization score $= D_{\text{region}} / D_{\text{global}}$, where
$D_{\text{region}}$ is direction instability computed on the masked gene
subset and $D_{\text{global}}$ is the standard 978-gene instability.

MOA prediction uses macro-averaged one-vs-rest AUROC: for each MOA class
with $\geq 10$ drugs, we compute the AUROC of the score (raw bracket or
localization score) as a predictor of class membership, then average
unweighted across classes.

\emph{Decision criterion (pre-registered):} Macro-averaged AUROC from
localization score exceeds that from raw bracket by $\geq 0.05$.
Bootstrap 95\% CI ($n = 2{,}000$ resamples) on the gap is reported.

\subsection{Results}

We computed localization scores for all drugs with annotated targets
and $\geq 5$ cell lines.  The macro-averaged AUROC gap is $+0.018$
(below the $+0.05$ threshold).  Per-MOA breakdown reveals that
localization improves prediction for intracellular enzyme inhibitors
but hurts for receptor-mediated drugs:

\begin{itemize}
\item Topoisomerase inhibitors: gap $= +0.49$
\item CDK inhibitors: gap $= +0.26$
\item Acetylcholine receptor modulators: gap $= -0.20$
\item PPAR agonists: gap $= -0.19$
\end{itemize}

\subsection{Discussion}

H4 is \textbf{not confirmed} at the pre-registered threshold.
The macro-averaged gap ($+0.018$) is positive but too small to
constitute evidence that gene-set restriction uniformly improves MOA
prediction.

The per-class pattern reveals a structural reason: localization helps
when the target's downstream transcriptional program is compact (e.g.,
topoisomerases directly affect a defined set of DNA-damage-response
genes), but hurts when the target operates through diffuse signaling
(e.g., acetylcholine receptors activate broad second-messenger
cascades whose transcriptomic footprint is distributed across many
genes, few of which appear in the top-100 shRNA mask).

This makes Rung~4 domain-dependent: the localization score adds value
for drugs with compact transcriptional targets but is actively
misleading for drugs with diffuse mechanisms.  The shRNA-derived region
mask is a further limitation: off-target effects in shRNA knockdowns
contaminate the consensus signature, potentially including irrelevant
genes and excluding relevant ones.

%% ===================================================================
\section{Dose-direction separability: an underpowered registered test}

A pre-registered dose-direction stability test (6{,}169
drug$\times$cell-line pairs) produced low within-pair cosines (median
0.19; 86.5\% below 0.5), failing both pre-registered criteria. Post-hoc
analysis revealed this failure is dominated by signature-level noise:
80\% of dose-level groups contain a single replicate, and
near-identical doses (ratio $< 1.01$) yield median cosine = 0.11,
establishing a noise floor close to the overall result. The test lacks
statistical power to distinguish genuine dose-dependent direction
changes from single-replicate noise in 978 dimensions. No a priori
power analysis preceded the pre-registration; replicate depth should
have been assessed before committing to specific thresholds.
Direction instability operates at a fixed dose across cell lines, a
regime where the relevant separability condition---that cell-line
response magnitude does not affect direction---is weaker than
dose-independence and holds trivially. The stronger claim (direction
independent of dose) remains empirically unresolved.

\section{HP2: Pathway function prediction from corrected Perturb-seq distances}
\label{sec:hp2}
%% ===================================================================

\subsection{Hypothesis}

If geometric distances between perturbation responses carry functional
information, then genes with similar GO (Gene Ontology) pathway
annotations should have more similar perturbation geometries than
functionally unrelated genes.  HP2 predicts a positive correlation
between pathway overlap (GO Jaccard similarity) and geometric similarity
(geodesic distance in the Grassmannian) for corrected Perturb-seq
distances.

\subsection{Methods}

For each pair of CRISPRi-perturbed genes in the Replogle et al.\ (2022)
Perturb-seq dataset, we computed the GO Jaccard similarity (intersection
over union of annotated GO biological process terms) and the pairwise
geodesic distance in the Grassmannian (from the 5-dimensional principal
subspace of each gene's perturbation response across cells).  Both raw
and essential-gene-corrected geodesic distances were tested.

Spearman correlation between GO Jaccard similarity and geodesic distance
was computed across all gene pairs with at least one shared GO term.

\emph{Decision criterion (pre-registered):} Spearman $|\rho| > 0.10$.

\subsection{Results}

GO Jaccard overlap has near-zero correlation with geodesic distance:
$|\rho| < 0.025$ for both raw and corrected forms.

\subsection{Discussion}

HP2 is a \textbf{null result}.  Geometric similarity of CRISPRi
perturbation responses does not predict functional pathway relatedness
as measured by GO annotation overlap.  Two non-exclusive explanations
apply.  First, GO annotations are incomplete and inconsistently granular:
two genes in the same pathway may share no GO terms if one is annotated
at a broad level (``regulation of transcription'') and the other at a
specific level (``positive regulation of NF-$\kappa$B signaling'').
Second, geometric similarity captures transcriptomic response similarity,
which may reflect shared downstream effects (convergent regulation) rather
than shared pathway membership.  Genes in different pathways can produce
similar transcriptomic perturbations if they converge on the same
effector genes.

This null establishes that geometric perturbation distance and functional
pathway annotation measure distinct biological properties.  Interpreting
geometric proximity as pathway co-membership would be unwarranted.

%% ===================================================================
\section{HDAC drug list for sensitivity analysis}
\label{sec:hdac_list}
%% ===================================================================

Table~\ref{tab:hdac_drugs} lists the 20 HDAC-targeting drugs removed in
the HDAC-removal sensitivity analysis (main text, Section~5.3).  Drugs
were identified by Drug Repurposing Hub target annotation matching
HDAC1, HDAC2, or HDAC6.  A synonym check for ``histone deacetylase,''
``class~I HDAC,'' ``pan-HDAC,'' and ``sirtuin'' confirmed no drugs were
missed by the primary filter.

\begin{table}[htbp]
\centering
\caption{HDAC-targeting drugs removed in the sensitivity analysis.
$D$ = raw direction instability.  All 20 drugs had $\geq 5$ cell lines.}
\label{tab:hdac_drugs}
\begin{tabular}{@{}lll@{}}
\toprule
Drug & HDAC target(s) & $D$ \\
\midrule
vorinostat      & HDAC1, HDAC2, HDAC6 & 0.490 \\
trichostatin-A  & HDAC1, HDAC2, HDAC6 & 0.550 \\
panobinostat    & HDAC1, HDAC2, HDAC6 & 0.534 \\
belinostat      & HDAC1, HDAC2, HDAC6 & 0.542 \\
PCI-24781       & HDAC1, HDAC2, HDAC6 & 0.529 \\
scriptaid       & HDAC1, HDAC2        & 0.631 \\
dacinostat      & HDAC1, HDAC2, HDAC6 & 0.572 \\
givinostat      & HDAC1, HDAC2        & 0.619 \\
mocetinostat    & HDAC1, HDAC2        & 0.651 \\
entinostat      & HDAC1               & 0.684 \\
tacedinaline    & HDAC1               & 0.712 \\
droxinostat     & HDAC6               & 0.734 \\
pyroxamide      & HDAC1               & 0.745 \\
tubastatin-A    & HDAC6               & 0.715 \\
PCI-34051       & HDAC8               & 0.889 \\
JNJ-26481585    & HDAC1, HDAC2, HDAC6 & 0.548 \\
Merck60         & HDAC1               & 0.692 \\
NSC-3852        & HDAC1               & 0.706 \\
phenylbutyrate  & HDAC1, HDAC2        & 0.768 \\
lovastatin      & HDAC2               & 0.913 \\
\bottomrule
\end{tabular}

\smallskip
\noindent\textit{Note:} Lovastatin is primarily an HMG-CoA reductase
inhibitor but carries a secondary HDAC2 annotation in the Drug
Repurposing Hub.  PCI-34051 targets HDAC8, which was not in the primary
filter but was included via the synonym check.
\end{table}

After removing all 20 drugs, the phenotype-projected correlation is
$\rho = 0.3759$ ($n = 775$, $p = 2.0 \times 10^{-27}$), compared with
$\rho = 0.3756$ on the full 795-drug set
($\Delta\rho = 3 \times 10^{-4}$).  The H3 result is not driven by the
HDAC drug class.

\end{document}
